%==================================================
\section*{녹색전환 분야}

\subsection*{1. 개요}

녹색전환은 기후위기 대응, 탄소중립, 자원순환 및 생태보전을 중심으로
농업ㆍ농촌의 생산체계와 자원관리 방식 전반의 구조적 변화를 의미한다.

최근 기후위기 심화, 식량안보 불확실성 확대, 국제 탄소규범 강화 등은
농업생산기반 관리기관으로서 한국농어촌공사의 역할 재정립을 요구하고 있다.

특히 공사의 핵심 기능인 농업용수 관리, 농지 정비, 배수개선,
농촌공간개발 기능은 녹색전환과 직접적으로 연계되는 영역으로,
기존의 시설 유지관리 중심 기능에서 벗어나
기후적응형 자원관리, 탄소저감형 기반시설 운영,
친환경 농촌공간 조성 기능으로 확대될 필요가 있다.

\subsection*{2. 주요 연구 내용}

본 연구의 녹색전환 분야는 다음의 네 단계로 구성한다.

\begin{enumerate}
    \item 녹색전환의 대내외 동향 및 농업ㆍ농촌 영향 분석
    \item 국내외 정책ㆍ제도 및 사례 검토
    \item 公社의 현행 기능 및 사업체계 진단
    \item 公社 기능 재정립 및 실천과제 도출
\end{enumerate}

%==================================================
\subsection*{3. 녹색전환의 대내외 동향 및 농업ㆍ농촌 영향 분석}

\subsubsection*{가. 주요 동향 분석}

최근 농업ㆍ농촌을 둘러싼 대내외 환경은 기후위기 심화, 국제 환경규범 강화,
식량안보 불확실성 확대, 자원순환 및 생태보전 요구 증대 등으로 급격한 구조변화를 겪고 있다.
이러한 변화는 농업생산기반의 안정성과 회복력을 저하시킬 수 있으며,
농업용수 관리, 농지관리, 재해 대응, 농촌공간 관리 등
한국농어촌공사의 핵심 기능 전반에 새로운 대응을 요구하고 있다.

\begin{itemize}

    \item \textbf{기후위기 심화 및 이상기후 확대}

    \begin{itemize}
        \item \textbf{활용 가능 자료}
        \begin{itemize}
            \item IPCC AR6 보고서
            \item 기상청 「2024년 이상기후 보고서」, 「우리나라 113년 기후변화 분석 보고서」
            \item 환경부 국가기후위기적응대책
            \item 통계청 자연재난 피해 통계
        \end{itemize}

        \item \textbf{확인할 내용}
        \begin{itemize}
            \item 평균기온 상승 추세
            \item 강수 패턴 변화
            \item 가뭄/홍수/폭염 발생 빈도 변화
            \item 작물 생산성 변화 및 재배적지 이동
            \item 농업용수 수급 영향
        \end{itemize}
    \end{itemize}

    \item \textbf{국제 탄소중립 및 환경규범 강화}

    \begin{itemize}
        \item \textbf{활용 가능 자료}
        \begin{itemize}
            \item UNFCCC / 국가 NDC 자료
            \item 2024 국가 온실가스 인벤토리 보고서
            \item EU Green Deal
            \item Farm to Fork Strategy
            \item 2050 농식품 탄소중립 추친전력
            \item 2050 탄소중립 녹색성장 기본계획
        \end{itemize}

        \item \textbf{확인할 내용}
        \begin{itemize}
            \item 농업부문 온실가스 배출 현황
            \item 국가별 감축 목표
            \item 저탄소 농업기술 및 정책
            \item 재생에너지 활용 확대 정책
        \end{itemize}
    \end{itemize}

    \item \textbf{식량안보 불확실성 확대}

    \begin{itemize}
        \item \textbf{활용 가능 자료}
        \begin{itemize}
            \item FAO Statistical Yearbook
            \item Global Report on Food Crises
            \item USDA 세계 곡물 수급 전망
            \item KREI 농업전망
        \end{itemize}

        \item \textbf{확인할 내용}
        \begin{itemize}
            \item 세계 식량가격 변동 추이
            \item 주요 곡물 수급 변화
            \item 국내 곡물자급률
            \item 공급망 리스크 및 수입 의존도
        \end{itemize}
    \end{itemize}

    \item \textbf{자원순환 및 생태보전 정책 강화}

    \begin{itemize}
        \item \textbf{활용 가능 자료}
        \begin{itemize}
            \item Strategic work of FAO for Sustainable Food and Agriculture
            \item OECD Circular Economy 관련 보고서
            \item 기후에너지환경부 자원순환기본계획
            \item 환경부 물 재이용 기본계획
            \item 국가생물다양성전략
        \end{itemize}

        \item \textbf{확인할 내용}
        \begin{itemize}
            \item 물 재이용 정책
            \item 토양ㆍ수질 관리 정책
            \item 생태계 기반 관리 사례
            \item 생물다양성 보전 정책
        \end{itemize}
    \end{itemize}

\end{itemize}

\subsubsection*{나. 농업ㆍ농촌 영향 분석}

녹색전환과 관련된 대내외 환경변화는 농업 생산체계, 자원관리 방식,
재해 대응체계, 농촌공간 구조 전반에 영향을 미치고 있다.
이에 따라 농업ㆍ농촌 부문에 미치는 영향을 분야별로 분석하고,
한국농어촌공사의 기능 변화와 연계할 필요가 있다.

\begin{itemize}

    \item \textbf{농업 생산성 변화 및 재배적지 이동}

    \begin{itemize}
        \item \textbf{활용 가능 자료}
        \begin{itemize}
            \item FAO 작물 생산량 및 생산성 통계
            \item 농촌진흥청 작물별 재배적지 변화 연구
            \item KREI 농업 생산성 및 식량수급 전망자료
            \item 국내 기후변화 시나리오 기반 작물 영향 연구
        \end{itemize}

        \item \textbf{확인할 내용}
        \begin{itemize}
            \item 주요 작물 생산량 변화 추이
            \item 단수 및 생산성 변화
            \item 재배 가능 지역 북상/축소 여부
            \item 품목 구조 변화 가능성
        \end{itemize}
    \end{itemize}

    \item \textbf{농업용수 수급 불안정 및 물관리 중요성 확대}

    \begin{itemize}
        \item \textbf{활용 가능 자료}
        \begin{itemize}
            \item 한국농어촌공사 농업용수 관련 통계
            \item 환경부 / 국토부 물관리 정책자료
            \item 기상청 강수량 및 가뭄 통계
            \item 국가 물관리 기본계획
            \item FAO Water-related agricultural reports
        \end{itemize}

        \item \textbf{확인할 내용}
        \begin{itemize}
            \item 농업용수 부족지역 현황
            \item 강수 패턴 변화와 용수 공급 영향
            \item 물 재이용 필요성
            \item 스마트 물관리 도입 필요성
        \end{itemize}
    \end{itemize}

    \item \textbf{재해위험 증가와 재해예방형 기반정비 수요 증가}

    \begin{itemize}
        \item \textbf{활용 가능 자료}
        \begin{itemize}
            \item 행정안전부 재해연보
            \item 기상청 자연재해 통계
            \item 한국농어촌공사 재해복구 사업자료
            \item 농림축산식품부 재해예방 정책자료
        \end{itemize}

        \item \textbf{확인할 내용}
        \begin{itemize}
            \item 농업재해 피해액 추이
            \item 침수 / 가뭄 / 염해 발생 현황
            \item 배수개선사업 수요 증가 여부
            \item 재해예방형 기반시설 투자 필요성
        \end{itemize}
    \end{itemize}

    \item \textbf{농촌공간의 친환경적 재구조화 필요성 확대}

    \begin{itemize}
        \item \textbf{활용 가능 자료}
        \begin{itemize}
            \item OECD Rural Policy / Rural 3.0
            \item 농림축산식품부 농촌공간계획 관련 자료
            \item 국토연구원 / KREI 농촌공간 재구조화 연구
            \item 환경부 생태복원 및 공간관리 정책자료
        \end{itemize}

        \item \textbf{확인할 내용}
        \begin{itemize}
            \item 농촌공간 이용구조 변화
            \item 친환경 공간조성 정책 방향
            \item 생활SOC 및 정주환경 개선 수요
            \item 생태친화형 농촌개발 방향
        \end{itemize}
    \end{itemize}

\end{itemize}

%==================================================
%==================================================
\subsection*{4. 국내외 정책ㆍ제도 및 사례 검토}

\subsubsection*{가. 국외 정책ㆍ제도 및 사례}

\begin{itemize}

    \item \textbf{EU 녹색전환 및 농업정책}

    \begin{itemize}
        \item \textbf{활용 가능 자료}
        \begin{itemize}
            \item European Green Deal
            \item Farm to Fork Strategy
            \item Common Agricultural Policy (CAP)
        \end{itemize}

        \item \textbf{확인할 내용}
        \begin{itemize}
            \item 저탄소 농업정책 방향
            \item 친환경 식량체계 구축 전략
            \item 농업환경지불 및 규제체계
        \end{itemize}
    \end{itemize}

    \item \textbf{미국 기후스마트농업 및 인프라 정책}

    \begin{itemize}
        \item \textbf{활용 가능 자료}
        \begin{itemize}
            \item USDA Climate-Smart Agriculture and Forestry
            \item Inflation Reduction Act (IRA)
            \item USDA Climate Hubs 자료
        \end{itemize}

        \item \textbf{확인할 내용}
        \begin{itemize}
            \item 기후스마트농업 지원정책
            \item 재생에너지 및 탄소감축 지원
            \item 농업재해 대응 및 보험제도
        \end{itemize}
    \end{itemize}

    \item \textbf{일본 기후적응형 농업 및 재해예방 정책}

    \begin{itemize}
        \item \textbf{활용 가능 자료}
        \begin{itemize}
            \item 일본 농림수산성(MAFF) 정책자료
            \item 국토강인화 기본계획
            \item 농업농촌정비사업 사례
        \end{itemize}

        \item \textbf{확인할 내용}
        \begin{itemize}
            \item 재해예방형 기반정비 정책
            \item 농업용수 및 배수관리 정책
            \item 농촌공간 재구조화 전략
        \end{itemize}
    \end{itemize}

    \item \textbf{네덜란드 스마트농업 및 물관리 정책}

    \begin{itemize}
        \item \textbf{활용 가능 자료}
        \begin{itemize}
            \item Dutch Delta Programme
            \item 스마트 물관리 및 스마트팜 정책자료
        \end{itemize}

        \item \textbf{확인할 내용}
        \begin{itemize}
            \item 물관리 거버넌스 체계
            \item 디지털 기반 농업관리 기술
            \item 기후적응형 농업 인프라 구축 전략
        \end{itemize}
    \end{itemize}

\end{itemize}

\subsubsection*{나. 국내 정책ㆍ제도 및 사례}

\begin{itemize}

    \item \textbf{탄소중립 및 기후위기 대응 정책}

    \begin{itemize}
        \item \textbf{활용 가능 자료}
        \begin{itemize}
            \item 2050 탄소중립 녹색성장 기본계획
            \item 국가 온실가스 감축목표(NDC)
            \item 국가기후위기적응대책
        \end{itemize}

        \item \textbf{확인할 내용}
        \begin{itemize}
            \item 국가 차원의 녹색전환 방향
            \item 농업부문 적응 및 완화 전략
        \end{itemize}
    \end{itemize}

    \item \textbf{농업ㆍ농촌 탄소중립 및 스마트농업 정책}

    \begin{itemize}
        \item \textbf{활용 가능 자료}
        \begin{itemize}
            \item 농업ㆍ농촌 탄소중립 추진전략
            \item 스마트농업 확산 정책
        \end{itemize}

        \item \textbf{확인할 내용}
        \begin{itemize}
            \item 저탄소 농업기술 확산 정책
            \item 디지털농업 및 생산성 향상 전략
        \end{itemize}
    \end{itemize}

    \item \textbf{농업용수 및 기반정비 정책}

    \begin{itemize}
        \item \textbf{활용 가능 자료}
        \begin{itemize}
            \item 농업용수 관리 정책자료
            \item 농업생산기반정비사업 정책자료
            \item 재해예방형 기반정비 관련 계획
        \end{itemize}

        \item \textbf{확인할 내용}
        \begin{itemize}
            \item 용수공급 안정화 정책
            \item 재해예방형 기반정비 방향
            \item 물관리 디지털화 방향
        \end{itemize}
    \end{itemize}

    \item \textbf{한국농어촌공사 추진 사례}

    \begin{itemize}
        \item \textbf{활용 가능 자료}
        \begin{itemize}
            \item 스마트 물관리 사례
            \item 농촌공간정비사업 사례
            \item 친환경 에너지 및 자원순환 사업 사례
        \end{itemize}

        \item \textbf{확인할 내용}
        \begin{itemize}
            \item 주요 사업 구조
            \item 적용 기술 및 제도
            \item 향후 확대 가능성
        \end{itemize}
    \end{itemize}

\end{itemize}

%==================================================
\subsection*{5. 公社의 현행 기능 및 사업체계 진단}

\begin{itemize}

    \item \textbf{활용 가능 자료}
    \begin{itemize}
        \item 한국농어촌공사 사업계획서
        \item 주요 사업 실적자료
        \item 내부 기능별 업무자료
    \end{itemize}

    \item \textbf{확인할 내용}
    \begin{itemize}
        \item 농업용수 관리 대응 수준
        \item 농지 정비 및 배수개선 대응 수준
        \item 농촌공간개발 대응 수준
        \item 기능별 한계와 보완 필요사항
    \end{itemize}

\end{itemize}

%==================================================
\subsection*{6. 公社 기능 재정립 및 실천과제 도출}

\begin{table}[h]
\centering
\renewcommand{\arraystretch}{1.3}
\begin{tabular}{p{5cm} p{9cm}}
\toprule
\rowcolor{gray!20}
현행 기능 & 녹색전환 대응 방향 \\
\midrule
농업용수 공급 & 기후적응형 스마트 물관리 \\
농지 정비 & 탄소저감형 농지관리 \\
배수개선 & 재해예방ㆍ회복형 기반시설 관리 \\
농촌공간개발 & 친환경ㆍ순환형 농촌공간 조성 \\
\bottomrule
\end{tabular}
\end{table}

\subsubsection*{단계별 실천과제}

\begin{itemize}
    \item 단기: 기후취약지역 진단 및 재해취약시설 보강
    \item 중기: 농업용수 재이용 확대 및 탄소중립형 사업 도입
    \item 장기: 기후적응형 농업 인프라 플랫폼 구축
\end{itemize}

%==================================================
\subsection*{7. 기대성과 및 활용방안}

\begin{itemize}
    \item 公社 기능 재정립 및 중장기 전략 수립 기초자료 활용
    \item 기존 사업 고도화 및 신규사업 발굴
    \item 정책 및 제도개선 검토자료 활용
\end{itemize}

%==================================================
\subsection*{8. 정리}

\begin{center}
\textbf{녹색전환 주요 동향 분석}
\\[0.5em]
$\Downarrow$
\\[0.5em]
\textbf{농업ㆍ농촌 영향 분석}
\\[0.5em]
$\Downarrow$
\\[0.5em]
\textbf{정책ㆍ제도 및 사례 검토}
\\[0.5em]
$\Downarrow$
\\[0.5em]
\textbf{公社 기능 진단}
\\[0.5em]
$\Downarrow$
\\[0.5em]
\textbf{기능 재정립 및 실천과제 도출}
\end{center}




\newpage

\subsection*{부록. 계획서 수정할 내용}
\section{국내ㆍ외 기후위기, 탄소중립 및 식량위기 동향과 전망 분석}

\begin{itemize}
    \item 기후위기와 탄소중립 전환은 농업ㆍ농촌의 생산체계, 자원관리 방식, 공간 기능을 동시에 변화시키는 핵심 요인으로 작용함.
    
    \item 지구온난화에 따른 기온 상승, 강수 패턴 변화, 이상기후 증가는 농업 생산성과 품질 저하, 병해충 증가, 농업용수 수급 불안정 등을 초래할 가능성이 큼.
    
    \item 우리나라 농업생산 체계는 온대 기후에 적합한 작부체계를 기반으로 발전해 왔으나, 향후 기후환경 변화에 따라 재배적지 이동, 작부체계 변화, 아열대 작물 도입 등이 확대될 것으로 예상됨.
    
    \item 탄소중립 사회로의 전환은 농업부문에 규제 부담과 새로운 기회를 동시에 제공함.
    \begin{itemize}
        \item 농약, 화학비료, 제초제, 화석연료 사용 등에 대한 제약 강화 가능
        \item 축산부문 항생제 사용, 사육두수, 가축분뇨 처리 등에 대한 관리 강화 가능
        \item 저탄소ㆍ친환경 농축산물 생산, 자원순환형 유통ㆍ소비체계, 재생에너지 기반 농촌공간 조성 기회 확대
    \end{itemize}

    \item 농업ㆍ농촌 공간은 국가 탄소중립 실현의 주요 거점으로 재평가될 필요가 있음.
    \begin{itemize}
        \item 농지, 산림, 습지, 수자원, 바이오매스 등 다양한 자연자원 보유
        \item 탄소흡수원 확대, 재생에너지 생산, 생물다양성 보전, 녹색산업 육성 가능
        \item 농촌공간이 생산공간에서 생태ㆍ에너지ㆍ바이오산업 기반공간으로 확장될 가능성
    \end{itemize}

    \item 기후위기는 글로벌 식량위기와도 밀접하게 연결됨.
    \begin{itemize}
        \item 가뭄, 홍수, 태풍, 한파, 산불 등 자연재해의 빈도와 강도 증가
        \item 물 부족, 토양 악화, 에너지 가격 상승, 국제 분쟁, 공급망 교란과 결합
        \item 주요 식량 생산지역의 생산 불안정성 확대
    \end{itemize}

    \item 세계 식량 수출 구조가 일부 국가에 집중되어 있어, 기후재해나 지정학적 충격 발생 시 국제 곡물가격 급등과 공급 불안이 발생할 가능성이 큼.

    \item 우리나라는 곡물자급률이 낮고 주요 곡물의 해외 의존도가 높기 때문에, 글로벌 식량위기의 심화는 국내 식량안보와 농업생산기반 안정성 문제로 직결됨.

    \item 따라서 녹색전환 시대의 농업ㆍ농촌 정책은 기후위기 대응, 탄소중립 이행, 식량안보 확보를 통합적으로 고려해야 하며, 농업생산기반은 기후재해 대응, 수자원 관리, 탄소저감, 생태보전, 식량안보를 함께 담당하는 전략적 기반으로 재정의될 필요가 있음.
\end{itemize}


\section{녹색전환 관점의 공사 역할 재정립 및 실천과제}

\begin{itemize}
    \item 녹색전환 시대의 한국농어촌공사는 기존의 농업생산기반 관리기관을 넘어 기후위기 대응, 탄소중립 이행, 식량안보 확보, 농촌자원 관리 기능을 함께 수행하는 공공기관으로 역할을 재정립할 필요가 있음.

    \item 농업기반 분야에서는 농업용수 공급, 농지정비, 배수개선 등 기존 기능을 기후적응형 생산기반 관리체계로 고도화해야 함.
    \begin{itemize}
        \item 가뭄ㆍ홍수 대응형 농업용수 관리체계 구축
        \item 저수지, 양배수장, 수로, 배수시설의 기후취약성 평가
        \item 재해위험 지역 중심의 농업기반시설 우선 보강
        \item 농업용수 재이용 및 순환형 물관리 체계 확대
    \end{itemize}

    \item 식량위기 대응 측면에서는 국내 농업생산기반의 유지와 보전 기능을 강화해야 함.
    \begin{itemize}
        \item 우량농지 보전 및 농지정비 기능 강화
        \item 전략작물 생산기반 지원
        \item 논의 다기능 활용 및 기후적응형 작부체계 전환 지원
        \item 안정적 용수공급과 배수개선을 통한 식량작물 생산 안정성 확보
    \end{itemize}

    \item 탄소중립 측면에서는 공사의 기존 사업을 저탄소ㆍ친환경 방식으로 전환할 필요가 있음.
    \begin{itemize}
        \item 농업기반시설의 에너지 효율화
        \item 저탄소 자재 및 친환경 공법 적용 확대
        \item 농지ㆍ저수지ㆍ수리시설을 활용한 재생에너지 사업 검토
        \item 바이오매스, 에너지작물 등 농촌자원 활용형 사업모델 발굴
    \end{itemize}

    \item 농촌개발 분야에서는 농촌공간을 탄소중립과 생태전환의 거점으로 전환하는 역할이 중요함.
    \begin{itemize}
        \item 농촌자원 데이터베이스 구축 및 모니터링 체계 마련
        \item 농촌 어메니티 자원 및 생물다양성 지표 개발
        \item 습지, 저수지, 수로, 하천 주변 수변공간의 체계적 관리
        \item 재생에너지 자립형 농촌공간 조성
        \item 녹색산업 및 농생명산업 기반 농촌개발모델 발굴
    \end{itemize}

    \item 공사의 녹색전환 대응전략은 다음 네 가지 방향으로 정리할 수 있음.
    \begin{enumerate}
        \item 기후위기 대응형 스마트 물관리 및 재해예방형 기반시설 관리체계 구축
        \item 식량위기 대응형 농업생산기반 유지 및 전략작물 생산기반 지원
        \item 탄소중립형 농업ㆍ농촌 인프라 확대
        \item 농촌자원의 생태적 가치를 활용한 친환경 농촌공간 관리체계 구축
    \end{enumerate}
\end{itemize}

\newpage

\subsection*{부록. 논문 및 자료 검색 키워드}

\subsubsection*{1. 주요 동향 분석}

\begin{itemize}

    \item \textbf{기후위기 심화 및 이상기후 확대}
    \begin{itemize}
        \item climate change agriculture
        \item climate risk agriculture
        \item extreme weather agriculture
        \item agricultural drought flood impact
        \item crop productivity climate change
        \item 기후변화 농업 영향
        \item 이상기후 농업 피해
    \end{itemize}

    \item \textbf{국제 탄소중립 및 환경규범 강화}
    \begin{itemize}
        \item carbon neutrality agriculture
        \item net zero agriculture
        \item agricultural greenhouse gas mitigation
        \item decarbonization agriculture
        \item low-carbon agriculture policy
        \item 탄소중립 농업
        \item 농업 온실가스 감축
    \end{itemize}

    \item \textbf{식량안보 불확실성 확대}
    \begin{itemize}
        \item food security climate change
        \item global food supply chain risk
        \item grain self-sufficiency
        \item food crisis agriculture
        \item 식량안보 기후변화
        \item 곡물자급률
    \end{itemize}

    \item \textbf{자원순환 및 생태보전 정책 강화}
    \begin{itemize}
        \item circular economy agriculture
        \item sustainable agriculture ecosystem
        \item biodiversity agriculture
        \item water reuse agriculture
        \item nature-based solutions agriculture
        \item 자원순환 농업
        \item 생태보전 농촌
    \end{itemize}

\end{itemize}

\subsubsection*{2. 농업ㆍ농촌 영향 분석}

\begin{itemize}

    \item \textbf{농업 생산성 변화 및 재배적지 이동}
    \begin{itemize}
        \item crop yield climate change
        \item crop suitability climate change
        \item agricultural productivity climate scenario
        \item 재배적지 변화
        \item 농업 생산성 변화
    \end{itemize}

    \item \textbf{농업용수 수급 불안정 및 물관리 중요성 확대}
    \begin{itemize}
        \item agricultural water management
        \item irrigation climate change
        \item drought irrigation management
        \item smart water management agriculture
        \item 농업용수 관리
        \item 스마트 물관리
    \end{itemize}

    \item \textbf{재해위험 증가와 재해예방형 기반정비 수요 증가}
    \begin{itemize}
        \item agricultural disaster management
        \item flood drainage agriculture
        \item climate resilience infrastructure agriculture
        \item 재해예방형 기반정비
        \item 농업재해 대응
    \end{itemize}

    \item \textbf{농촌공간의 친환경적 재구조화 필요성 확대}
    \begin{itemize}
        \item rural restructuring sustainability
        \item rural spatial planning
        \item eco-friendly rural development
        \item 농촌공간 재구조화
        \item 친환경 농촌개발
    \end{itemize}

\end{itemize}

\subsubsection*{3. 정책ㆍ제도 및 사례 검토}

\begin{itemize}

    \item \textbf{국제 논의 및 정책 동향}
    \begin{itemize}
        \item climate-smart agriculture policy
        \item agricultural adaptation policy
        \item sustainable food system policy
        \item agricultural policy climate change
    \end{itemize}

    \item \textbf{국내 정책 동향}
    \begin{itemize}
        \item Korea carbon neutrality agriculture
        \item Korea agricultural climate adaptation
        \item 한국 농업 탄소중립 정책
        \item 한국 농업 기후위기 적응
    \end{itemize}

    \item \textbf{국내외 사례 조사}
    \begin{itemize}
        \item smart irrigation case study
        \item climate-smart agriculture case study
        \item rural resilience case study
        \item 스마트 물관리 사례
        \item 기후스마트농업 사례
    \end{itemize}

\end{itemize}

\subsubsection*{4. 公社 기능 및 사업체계 진단}

\begin{itemize}
    \item public agency role climate adaptation
    \item agricultural infrastructure management
    \item public water management system
    \item 농업기반시설 관리
    \item 공공기관 기능 재편
\end{itemize}

\subsubsection*{5. 기능 재정립 및 실천과제 도출}

\begin{itemize}
    \item strategic roadmap public agency
    \item agricultural policy roadmap
    \item sustainable rural development strategy
    \item 공공기관 발전전략
    \item 농업정책 로드맵
\end{itemize}